\chapter{Geometry Generates the Gauge}
\label{chap:geometry-generates-gauge}

The last chapter named a value of the invariant the electron and was careful to call the naming a
naming. This chapter asks where that value comes from, and the answer turns out to be structural rather
than stipulated: the discriminating configuration whose second variation reads the electron is not put
in by hand but \emph{generated}, from geometry, by a procedure the construction builds and runs. The
result is that a gauge field --- the carrier physics pairs with a charge --- emerges from a geometric
program rather than being postulated alongside it. Before any of that, the fence: the construction is
unusually explicit, in its own notes, about the claim ceiling here. Its geometric program is
\emph{GR-like}, a discrete shadow of geometry and not General Relativity; what it generates is a
\emph{QED-like} program, a model with the shape of quantum electrodynamics and not the theory itself.
The construction's own phrase is the honest one to keep in view: a shadow, no more. What follows is the
generation of one shadow from another, which is a real and surprising thing, and is exactly that much.

\section{The geometric program}
\label{se:the-geometric-program}

The construction packages what it has built into a single object it calls a geometric program: the
labeled universe kernel of the sixteenth chapter, together with its Hilbert door and its anchor. Calling
this geometry earns a careful justification, and the construction gives one. A stiffness form on a graph
of nodes, with its couplings between neighbors and its anchor at the boundary, is a discrete stand-in
for a metric --- a record of how the configuration is bent and strained from place to place --- and the
invariant it carries is a discrete stand-in for a curvature, the second-order quantity a metric induces.
That is what earns the word geometric: not that the construction has built spacetime, but that its
finite operator has the shape of one, the same way its sample had the shape of a real number. The
program's invariant still detects zero, and its source --- the boundary where the bulk's accounts are
settled --- is the white-hole source met before, the place a charge's influence originates in the
record.

It is worth being concrete about the shadow, because the word geometric is doing real work and should
not do more than it can. A metric, in the continuous theory, tells you the distance between nearby
points, and from how that distance varies you read curvature --- the failure of the geometry to be flat,
the quantity that bends paths and sources gravity. The construction's stiffness form is the discrete
echo of exactly this. Its couplings between neighboring nodes say how far apart, in the energy's
reckoning, adjacent parts of the configuration sit; the variation of those couplings from node to node
is a discrete curvature, the second-order datum the form induces; and the anchor at the boundary is the
condition that pins the whole down, the way a boundary condition pins a metric. None of this is
spacetime. It is a finite graph with weights, and the claim is only that the weights are arranged in the
pattern a geometry would leave --- that the form has, in the precise structural sense the book has used
since it built a sample with the shape of a real number, the shape of a curved space. The construction's
gate for the claim is set exactly there and audited to stay there: geometric in shape, a shadow, no
more.

One honest dependency belongs in the open. The geometric side of the program rests on the Hilbert door
of the sixteenth chapter --- the completion the construction marks but does not build --- and so it
carries the one non-finite import the book has admitted, the limit imported at the threshold rather than
assembled. The construction notes this plainly: the geometric chair leans on that imported completion,
while the gauge program it will generate is constructive throughout, built and runnable without reaching
past the door. So even within this chapter there are two registers, and the construction keeps them
apart: a geometric program that borrows the completion at the threshold, and a generated gauge program
that borrows nothing.

This two-register structure is not a blemish to be hidden but the most honest thing in the chapter, and
it is worth seeing why the construction tolerates the import on the geometric side. The Hilbert door is
the completion the sixteenth chapter flagged: the step from the finite, refinable tower to the completed
space it approaches, the one place the construction reaches for a limit it cannot finitely assemble. The
geometric program wants that completed space, because geometry in the usual sense lives on a continuum,
and so the geometric side leans on the door. But --- and this is the point of keeping the registers
apart --- the gauge program generated from it does not. The boundary generator, the discriminating
action, the electron it produces, the holonomy it assigns: all of these are finite, constructive, and
runnable, and none of them reaches back through the door. So the import is quarantined. It sits on the
geometric input, declared and bounded, and does not leak into the gauge output, which is exactly the
separation a careful account of a half-finite, half-imported construction should maintain.

\section{The boundary generates the gauge}
\label{se:the-boundary-generates-the-gauge}

Now the chapter's engine. The construction builds a boundary procedure --- a generator that reads the
geometric program at its boundary and produces, as output, the discriminating action whose second
variation is the electron. The geometry is the input; the gauge action is the output; and the procedure
that turns one into the other lives at the boundary, the white-hole source where the geometry's accounts
are settled. The construction proves the generation works: the procedure run on the three-node geometric
program generates exactly the discriminating action, and that generated action produces the electron of
the previous chapter. The gauge is not postulated beside the geometry; it is funged out of the
geometry's boundary, assembled from the bending of the configuration into the action that reads a charge.

This is the picture the experiment under Aharonov and Bohm draws in the laboratory. A charged particle
carried around a loop picks up a phase --- a holonomy --- set by the gauge potential, and it does so even
where the field itself vanishes, so that the gauge effect is read off the geometry of the path rather
than off any local force. Geometry, in that experiment, generates gauge: the loop and the potential it
encloses determine a phase that is a genuine, measurable physical fact with no local field to account
for it. The construction's boundary generator is the same move made constructive. It reads the geometry
--- the loop, the bending, the boundary --- and produces the gauge action whose holonomy a charge would
feel; the phase the Aharonov--Bohm particle acquires is, in the construction, the reading the generated
action returns. The gauge tanges the loop: it selects, out of the geometry, exactly the holonomy that a
charge transported around it accumulates, and reports it as the action's content.

The worked picture is the cleanest the chapter has. Imagine the charge carried slowly around a closed
loop in the geometry, returning to where it started. In the continuous experiment it comes back not
quite as it left: it has acquired a phase, an angle, set entirely by what the loop encloses, even if the
charge never passed through any region where a force acted on it. The phase is a holonomy --- a quantity
that lives on the loop as a whole and cannot be read off any single point of it. The construction's
boundary generator produces exactly the action that assigns this holonomy. Hand it the geometry of the
loop --- the bending, the enclosed source at the white-hole boundary --- and it returns the gauge action
whose reading on a charge transported around the loop is precisely that accumulated phase. So the gauge
field is not a second thing laid over the geometry with its own independent existence; it is the
geometry's holonomy, the loop-wise residue of the bending, made into an action the construction can run.
Geometry on the input side, gauge on the output side, and a boundary procedure, proved correct, turning
the one into the other.

\section{The generated program}
\label{se:the-generated-program}

Run the generation to its end and the construction has produced not just an action but a whole program
--- the electron together with the gauge field that couples to it, a generated, internally certified
account with the shape of quantum electrodynamics. From the geometric program, by the boundary
generator, comes a gauge program: a charged object and the field it interacts through, both built, both
auditable, both following from the geometry rather than added to it. The construction's note for this is
the one Yang and Mills would recognize --- a gauge field arising from the demand of admissibility, the
requirement that the construction's bookkeeping stay consistent under the relabelings the fifth chapter
called gauge freedom. The experiment under the Yang--Mills name marks the sector: the gauge field is what
admissibility forces once a charge is present, the carrier that must exist for the labeled invariant to
stay consistent across the local frames it lives in.

The logic by which admissibility forces a gauge field is worth making explicit, because it is the same
logic the fifth chapter introduced and the whole book has relied on. A charge is a labeled invariant,
and the label is local: it is read in a frame, and different parts of the configuration may be labeled in
frames that do not agree. For the invariant to be a single consistent object across the whole
configuration, there must be a rule for carrying a label from one frame to its neighbor --- a connection,
in the geometric word --- and that rule, the thing that reconciles the local frames, is the gauge field.
It is not optional. Once a charge is present and its label is local, admissibility --- the demand that
the bookkeeping be consistent everywhere at once --- forces a connection to exist, and the connection is
the gauge potential whose holonomy the previous section described. This is what it means to say the
gauge is generated rather than postulated: the construction does not add a field by hand and check that
it is consistent; it requires consistency and finds that a field must be there, the carrier without
which the labeled invariant could not be single-valued across its frames.

And here the fence must be held hardest, because the temptation to overread is greatest. What the
construction has generated is a program with the \emph{structure} of quantum electrodynamics: a charged
particle, a gauge field, a coupling between them, all emerging from a geometric program with the
structure of general relativity. That two such shadows --- a discrete geometry and the gauge field it
forces --- stand in exactly the generative relation that geometry and gauge are believed to stand in for
the real theories is a striking structural fact, and the construction proves it. But it is a fact about
shadows. The generated program has no particular masses, no measured coupling constant, no claim to be
the electrodynamics of the world; it is QED-shaped, generated from GR-shaped geometry, and the word
\emph{shaped} is carrying the whole modeling claim. The construction shows that the shape of
electromagnetism can be made to fall out of the shape of geometry, which is a great deal, and stops
precisely there.

\section{What is demonstrated, and what is assumed}
\label{se:demonstrated-assumed-ch20}

The separation here has three layers, and they are worth keeping distinct.

What is demonstrated is the generation. The geometric program is built --- a kernel, an anchor, an
invariant that detects zero. The boundary generator is proved to produce the discriminating action from
that program, and the generated action is proved to produce the electron. And the whole is proved to
constitute a generated gauge program, a QED-shaped account following from the geometry. The gauge program
is constructive and runnable; none of the generation reaches past what the construction builds.

\begin{quote}
\emph{A note on the labels --- two shadows and one door.} Three things here are bridges of different
weights. That the finite stiffness program is \emph{geometry} is the modeling claim of the sixteenth
chapter carried forward: it is GR-shaped, a discrete shadow of a metric and its curvature, marked by the
construction's own gate as a shadow and no more. That what is generated is \emph{a gauge field} and a
\emph{QED program} is the second shadow: it is QED-shaped, the structure of electrodynamics generated
from the structure of geometry, with none of the physical theory's measured content. And the geometric
side's reliance on the imported Hilbert completion is the one non-finite assumption admitted, a door
flagged not built. The genuine, proved, surprising fact in the middle of all this hedging is the
generative relation itself: that the shape of gauge falls out of the shape of geometry by a boundary
procedure the construction runs. That relation is demonstrated; the identification of either shape with
the physics it resembles is named.
\end{quote}

What is assumed, then, is the two modeling identifications and the imported completion, on top of the
standing oracle --- and the construction is scrupulous that the result in the middle, the generation,
needs none of them to be true as a statement about its own shadows. With a gauge generated from geometry
and a charged particle falling out of the boundary, the construction has the makings of a field theory.
The next chapter follows the split between the geometric and gauge sectors to its first physical
consequence: that the mixed terms vanish in the interior, so that radiation, when it comes, is a
phenomenon of the boundary.
